\section{The Fusion Design Philosophy}

A system that allocates intelligence cannot be organized primarily around a
catalog of models. Models change too quickly, their strengths are conditional,
and a model name says little about what a particular task needs next. Fusion is
instead organized around the information state of the task.

\principle{Organize computation around uncertainty, not models.}

This is the foundational design principle of Fusion. The system begins with
what is unresolved: a missing fact, an untested assumption, a fragile plan, an
ambiguous preference, an incomplete search, or an unobserved state of the
world. Only then does it ask which computation could resolve that uncertainty.

The principle makes the architecture durable. Models, providers, and tools can
change without changing the object the system is trying to control. A new
capability enters Fusion by demonstrating what uncertainty it can resolve, in
which role, under which conditions. The orchestration layer does not need a
permanent ranking of models; it needs an evolving understanding of their
conditional value.

\subsection{From confidence to decision-relevant uncertainty}

Language models contain useful internal signals about uncertainty. Token
entropy captures uncertainty over expression, while semantic entropy can
better reflect disagreement over meaning
\citep{kuhn2023semantic,farquhar2024semantic}. A compound system must go one
step further. It must care about \emph{decision-relevant uncertainty}: what is
unknown that could change the answer, action, or confidence required to
proceed?

This distinction prevents activity from masquerading as progress. Several
models can agree because they share the same training patterns. Many diverse
answers can be irrelevant to the load-bearing decision. Conversely, one small
piece of external evidence can collapse an entire branch of reasoning.

Fusion therefore treats each potential call as an information acquisition
decision. In practical terms, its value is:

\begin{equation}
  \text{value of a call}
  = \text{expected decision improvement}
  - \text{compute} - \text{latency} - \text{risk}.
\end{equation}

The system does not need a perfect Bayesian model of an open-ended task for
this principle to be useful. It needs structured approximations: what claims
are unresolved, what evidence already exists, which failure modes matter, how
ready the workflow is to act, and which source has historically changed similar
decisions. These approximations can be inspected today and learned from
outcomes over time.

\subsection{Models as conditional sources of evidence}

Under this philosophy, a model is not a scalar score and a call is not a vote.
A model is a conditional source of evidence. Its contribution depends on the
task, role, context, evidence already present, provider surface, and budget.
The same model can be highly valuable as a critic and redundant as a second
solver. A lower-cost model can have high marginal value when it searches a
different region or detects a different class of error. A frontier model can
be indispensable when integration, ambiguity, or stakes demand stronger
judgment.

This view gives cost its proper technical meaning. Cost constrains which
observations can be acquired and how often a system can learn from them. A
model's practical capability is inseparable from the conditions under which it
can be deployed. Understanding cost, latency, reliability, complementarity,
and error correlation is part of understanding the model itself.

\subsection{Plural intelligence, singular authority}

Multiple perspectives create value only if the system remains coherent.
Fusion therefore pairs its information principle with an authority principle:

\principle{One state, many perspectives, one commit.}

The system owns one canonical representation of the task and its history.
Models and tools receive bounded views and return attributed evidence. They do
not independently redefine the task, mutate the shared state, or issue
client-visible actions. One authoritative path integrates the evidence and
commits the next answer or transition.

This separates epistemic plurality from operational plurality. Fusion can
search widely without becoming a swarm of conflicting agents. It can challenge
a plan without creating competing world states. And it can remain compatible
with the simple external contract of one model, even while drawing on a much
richer internal supply of intelligence.

Together, the two principles define Fusion's design philosophy:

\begin{itemize}
  \item acquire computation only against a decision-relevant uncertainty;
  \item interpret outputs as scoped, attributed evidence rather than authority;
  \item preserve one canonical task state across the full workflow; and
  \item stop when further evidence is unlikely to change the commitment.
\end{itemize}
